%% LyX 2.4.0 created this file.  For more info, see https://www.lyx.org/.
%% Do not edit unless you really know what you are doing.
\documentclass[spanish]{article}
\usepackage[T1]{fontenc}
\usepackage[utf8]{inputenc}
\usepackage{amsmath}
\usepackage{amssymb}
\usepackage{babel}
\addto\shorthandsspanish{\spanishdeactivate{~<>.}}
\deactivatequoting

\begin{document}
\title{Exploring LFV Higgs decays in the Three Higgs Doublet Model}
\author{M. Zeleny-Mora$^{1}$\thanks{moiseszeleny@gmail.com}, M. Mondragón$^{1}$\thanks{myriam@fisica.unam.mx},
T. A. Valencia-Pérez$^{1}$\thanks{tvalencia@fisica.unam.mx}\\
$^{1}$Instituto de Física, Universidad Nacional Autónoma de México,
Apdo. Postal 20-364, CDMX 01000, México.}
\maketitle
\begin{abstract}
In this work we consider the Three Higgs Doublet model with $S_{3}$
symmetry, we focus on the lepton sector and explore lepton flavor
violating Higgs decays.
\end{abstract}

\section{Introduction}

In the search of new

\section{The model}

In this work we consider the Three Higgs Doublet Model with $S_{3}$
symmetry (3HDMS3) \cite{Das:2015sca}. The discrete symmetry $S_{3}$
has three irreducible representations denoted by $\mathbf{1}$, $\mathbf{1^{\prime}}$
and $\mathbf{2}.$ In the basis where the generators of $S_{3}$group
in the $\mathbf{2}$ representation is given by
\[
a=\begin{bmatrix}-\frac{1}{2} & \frac{\sqrt{3}}{2}\\
-\frac{\sqrt{3}}{2} & -\frac{1}{2}
\end{bmatrix},\qquad b=\begin{bmatrix}\frac{1}{2} & \frac{\sqrt{3}}{2}\\
\frac{\sqrt{3}}{2} & -\frac{1}{2}
\end{bmatrix}
\]
 where $a$ is of order 3 and $b$ is of order 2. In the lepton sector,
the fields are assigned the following representations of the $S_{3}$:
\begin{align}
\mathbf{2:} & \begin{bmatrix}L_{1}\\
L_{2}
\end{bmatrix},\,\begin{bmatrix}\nu_{1R}\\
\nu_{2R}
\end{bmatrix},\,\begin{bmatrix}e_{1R}\\
e_{2R}
\end{bmatrix}\nonumber \\
\mathbf{1}: & L_{S},\nu_{SR},e_{SR},\label{eq:S3_representations_leptons}
\end{align}
where the $L_{i}$ with $i=1,2,3$, are the usual $SU\left(2\right)_{L}$
lepton doublets, whereas the $\nu_{iR}$ and $e_{iR}$ are the right-handed
neutrinos and charged leptons fields which are singlets of the $SU\left(2\right)_{L}$
part of gauge symmetry. The square brackets in (\ref{eq:S3_representations_leptons})
represent the doublet representation under $S_{3}$, different to
the doublet representation under $SU\left(2\right)_{L}$. In the scalar
sector, there are three $SU\left(2\right)_{L}$ doublets $\phi_{i}$
and this transform under the $S_{3}$ symmetry as follows:
\[
\mathbf{2}:\begin{bmatrix}H_{1}\\
H_{2}
\end{bmatrix}\equiv\Phi,\qquad\mathbf{1}:H_{S}.
\]
where
\[
H_{1}=\frac{1}{\sqrt{2}}\begin{pmatrix}\phi_{1}+i\phi_{4}\\
\phi_{7}+i\phi_{10}
\end{pmatrix};\quad H_{2}=\frac{1}{\sqrt{2}}\begin{pmatrix}\phi_{2}+i\phi_{5}\\
\phi_{8}+i\phi_{11}
\end{pmatrix};
\]
\[
H_{S}=\frac{1}{\sqrt{2}}\begin{pmatrix}\phi_{3}+i\phi_{6}\\
\phi_{9}+i\phi_{12}
\end{pmatrix}.
\]
In the case of the most general scalar potential in the (3HDMS3),
considering that each $\phi_{i}$ ($i=7,8,9$) has a real vacuum expectation
value (VEV) $v_{i}$ non-zero. The minimization of the potential yields
to the relation
\begin{equation}
v_{1}=\sqrt{3}v_{2}.\label{eq:vev12_cond}
\end{equation}

The condition (\ref{eq:vev12_cond}), allows to $Z_{2}$ symmetry
survive the spontaneous symmetry breaking, or equivalently the VEVs
are unaffected by the transformation
\[
\Phi\to b\Phi.
\]
In this model, $v=\sqrt{v_{1}^{2}+v_{2}^{2}+v_{3}^{2}}=246$ GeV appears
in the masses of $W$ and $Z$ bosons. The VEV's be rewritten in spherical
coordinates as follows: 
\[
v_{1}=v\cos\varphi\sin\theta_{v},\quad v_{2}=v\sin\varphi\sin\theta_{v},\quad v_{3}=v\cos\theta_{v},
\]
in this form, the relations among vev's are parametrized in terms
of two angles
\[
\tan\varphi=\frac{v_{2}}{v_{1}},\quad\tan\theta_{v}=\frac{v_{2}}{v_{3}\sin\varphi}.
\]
Moreover, the value of $\varphi$ is fixed by the minimum of the potential
(\ref{eq:vev12_cond}), as follows: 
\begin{align*}
\tan\varphi= & \frac{1}{\sqrt{3}}\Rightarrow\tan\theta_{v}=\frac{2v_{2}}{v_{3}}\\
\sin\theta_{v}= & \frac{2v_{2}}{v}\Rightarrow\cos\theta_{v}=\frac{v_{3}}{v}.
\end{align*}

The CP even Higgs in the mass basis $h_{0}$, $H_{1}$ and $H_{2}$
are rotated as follows:
\[
\begin{pmatrix}\phi_{7}\\
\phi_{8}\\
\phi_{9}
\end{pmatrix}=R_{S}\begin{pmatrix}h_{1}\\
h_{0}\\
h_{2}
\end{pmatrix}
\]
where
\begin{equation}
R_{S}=\left(\begin{array}{ccc}
\sin\alpha\cos\varphi & -\sin\varphi & -\cos\alpha\cos\varphi\\
\sin\alpha\sin\varphi & \cos\varphi & -\cos\alpha\sin\varphi\\
\cos\alpha & 0 & \sin\alpha
\end{array}\right).\label{eq:RS}
\end{equation}

To straightforward comparison with the SM like we consider the Higgs
basis where only one Higgs doublet acquire vev. 
\[
\begin{pmatrix}\Phi_{\text{vev}}\\
\Phi_{1}\\
\Phi_{2}
\end{pmatrix}=R_{A}^{\top}\begin{pmatrix}H_{1}\\
H_{2}\\
H_{S}
\end{pmatrix},
\]
\begin{align}
R_{A}^{\top} & =\left(\begin{array}{ccc}
\sin\theta_{v}\cos\varphi & \sin\theta_{v}\sin\varphi & \cos\theta_{v}\\
-\sin\varphi & \cos\varphi & 0\\
-\cos\theta_{v}\cos\varphi & -\cos\theta_{v}\sin\varphi & \sin\theta_{v}
\end{array}\right),\label{eq:RAT}\\
 & =\left(\begin{array}{ccc}
\frac{\sqrt{3}v_{2}}{v} & \frac{v_{2}}{v} & \frac{v_{3}}{v}\\
-\frac{1}{2} & \frac{\sqrt{3}}{2} & 0\\
-\frac{\sqrt{3}v_{3}}{2v} & -\frac{v_{3}}{2v} & \frac{2v_{2}}{v}
\end{array}\right)
\end{align}
where
\[
\begin{aligned}\Phi_{\text{vev}} & =\binom{G^{\pm}}{\frac{1}{\sqrt{2}}\left(v+\widetilde{H}+iG_{0}\right)},\quad\Phi_{1}=\binom{H_{1}^{\pm}}{\frac{1}{\sqrt{2}}\left(\widetilde{H}_{a}+iA_{1}\right)},\\
\Phi_{2} & =\binom{H_{2}^{\pm}}{\frac{1}{\sqrt{2}}\left(\widetilde{H}_{b}+iA_{2}\right)}.
\end{aligned}
\]
In this basis, the Goldstone bosons $G_{0}$, $G^{\pm}$, pseudoscalars
$A_{1,2}$ and charged scalars $H_{1,2}^{\pm}$ appears explicitly.
Whereas the neutral part of the Higgs doublets, denoted by $\widetilde{H}$,
$\widetilde{H}_{a}$ and $\widetilde{H}_{b}$ do not correspond to
their mass eigenstates, but they are in the Higgs basis. So, we rotate
into the physical basis from the Higgs basis as follows
\[
\left(\begin{array}{c}
\widetilde{H}\\
\widetilde{H}_{a}\\
\widetilde{H}_{b}
\end{array}\right)=R_{H}\left(\begin{array}{l}
h_{1}\\
h_{0}\\
h_{2}
\end{array}\right),
\]
\begin{equation}
R_{H}=\left(\begin{array}{ccc}
\cos(\alpha-\theta_{v}) & 0 & \sin(\alpha-\theta_{v})\\
0 & 1 & 0\\
-\sin(\alpha-\theta_{v}) & 0 & \cos(\alpha-\theta_{v})
\end{array}\right).\label{eq:R_H}
\end{equation}
In consequence, we have
\begin{equation}
R_{S}=R_{A}R_{H}.\label{eq:RS_AH}
\end{equation}

and three scenarios:
\begin{enumerate}
\item Scenario A: $h_{1}$ SM-like In this case, 
\begin{align*}
h_{1} & =\widetilde{H}\Rightarrow\cos(\alpha-\theta_{v})=1\\
\alpha & =\theta_{v}\Rightarrow R_{H}=\mathbb{I},\Rightarrow R_{S}=R_{A}
\end{align*}
\begin{align}
\phi_{7}= & \frac{\sqrt{3}v_{2}}{v}h_{1}-\frac{1}{2}h_{0}-\frac{\sqrt{3}v_{3}}{2v}h_{2},\nonumber \\
\phi_{8}= & \frac{v_{2}}{v}h_{1}+\frac{\sqrt{3}}{2}h_{0}-\frac{v_{3}}{2v}h_{2},\nonumber \\
\phi_{9}= & \frac{v_{3}}{v}h_{1}+\frac{2v_{2}}{v}h_{2}\label{eq:phi789_h1}
\end{align}
\item Scenario B: $h_{2}$ SM-like, In this case
\begin{align*}
h_{2} & =\widetilde{H}\Rightarrow\sin(\alpha-\theta_{v})=1\\
\alpha & -\theta_{v}=\frac{\pi}{2}\\
\Rightarrow & R_{H}=\left(\begin{array}{ccc}
0 & 0 & 1\\
0 & 1 & 0\\
-1 & 0 & 0
\end{array}\right)\\
\Rightarrow & R_{S}=\begin{pmatrix}\frac{\sqrt{3}\cos\left(\theta_{v}\right)}{2} & -\frac{1}{2} & \frac{\sqrt{3}\sin\left(\theta_{v}\right)}{2}\\
\frac{\cos\left(\theta_{v}\right)}{2} & \frac{\sqrt{3}}{2} & \frac{\sin\left(\theta_{v}\right)}{2}\\
-\sin\left(\theta_{v}\right) & 0 & \cos\left(\theta_{v}\right)
\end{pmatrix}\\
\Rightarrow & R_{S}=\begin{pmatrix}\frac{\sqrt{3}v_{3}}{2v} & -\frac{1}{2} & \frac{\sqrt{3}v_{2}}{v}\\
\frac{v_{3}}{2v} & \frac{\sqrt{3}}{2} & \frac{v_{2}}{v}\\
-\frac{2v_{2}}{v} & 0 & \frac{v_{3}}{v}
\end{pmatrix}
\end{align*}
\begin{align}
\phi_{7}= & \frac{\sqrt{3}v_{3}}{2v}h_{1}-\frac{1}{2}h_{0}+\frac{\sqrt{3}v_{2}}{v}h_{2},\nonumber \\
\phi_{8}= & \frac{v_{3}}{2v}h_{1}+\frac{\sqrt{3}}{2}h_{0}+\frac{v_{2}}{v}h_{2},\nonumber \\
\phi_{9}= & -\frac{2v_{2}}{v}h_{1}+\frac{v_{3}}{v}h_{2}\label{eq:phi789_h2}
\end{align}
\item Scenario C: General $\alpha-\theta_{v}=\delta$
\begin{align*}
\alpha-\theta_{v} & =\delta\ll1\Rightarrow R_{H}\approx\left(\begin{array}{ccc}
\cos\left(\delta\right) & 0 & \sin\left(\delta\right)\\
0 & 1 & 0\\
-\sin\left(\delta\right) & 0 & \cos\left(\delta\right)
\end{array}\right).\\
\Rightarrow & R_{S}=\begin{pmatrix}\frac{\sqrt{3}v_{2}\cos\left(\delta\right)}{v}+\frac{\sqrt{3}v_{3}\sin\left(\delta\right)}{2v} & -\frac{1}{2} & \frac{\sqrt{3}v_{2}\sin\left(\delta\right)}{v}-\frac{\sqrt{3}v_{3}\cos\left(\delta\right)}{2v}\\
\frac{v_{2}\cos\left(\delta\right)}{v}+\frac{v_{3}\sin\left(\delta\right)}{2v} & \frac{\sqrt{3}}{2} & \frac{v_{2}\sin\left(\delta\right)}{v}-\frac{v_{3}\cos\left(\delta\right)}{2v}\\
-\frac{2v_{2}\sin\left(\delta\right)}{v}+\frac{v_{3}\cos\left(\delta\right)}{v} & 0 & \frac{2v_{2}\cos\left(\delta\right)}{v}+\frac{v_{3}\sin\left(\delta\right)}{v}
\end{pmatrix}
\end{align*}
\begin{align}
\phi_{7}= & \frac{\sqrt{3}}{2v}\left(v_{3}\sin\left(\delta\right)+2v_{2}\cos\left(\delta\right)\right)h_{1}-\frac{1}{2}h_{0}+\frac{\sqrt{3}}{2v}\left(2v_{2}\sin\left(\delta\right)-v_{3}\cos\left(\delta\right)\right)h_{2},\nonumber \\
\phi_{8}= & \frac{1}{2v}\left(v_{3}\sin\left(\delta\right)+2v_{2}\cos\left(\delta\right)\right)h_{1}+\frac{\sqrt{3}}{2}h_{0}+\frac{1}{2v}\left(2v_{2}\sin\left(\delta\right)-v_{3}\cos\left(\delta\right)\right)h_{2},\nonumber \\
\phi_{9}= & -\frac{2v_{2}\sin\left(\delta\right)-v_{3}\cos\left(\delta\right)}{v}h_{1}+\frac{1}{v}\left(v_{3}\sin\left(\delta\right)+2v_{2}\cos\left(\delta\right)\right)h_{2}.\label{eq:phi789_near_alignment}
\end{align}
\end{enumerate}

\subsection{Yukawa sector}

In the Yukawa sector for charged leptons 
\begin{align*}
\mathcal{L}_{Y_{\ell}}= & -\frac{1}{\sqrt{2}}Y_{1}^{\ell}\left[\bar{L}_{1}\ell_{1R}+\bar{L}_{2}\ell_{2R}\right]H_{S}-Y_{3}^{\ell}\bar{L}_{S}H_{S}\ell_{SR}\\
 & -\frac{1}{\sqrt{2}}Y_{2}^{\ell}\left[\left(\bar{L}_{1}H_{2}+\bar{L}_{2}H_{1}\right)\ell_{1R}+\left(\bar{L}_{1}H_{1}-\bar{L}_{2}H_{2}\right)\ell_{2R}\right]\\
 & -\frac{1}{\sqrt{2}}Y_{4}^{\ell}\bar{L}_{S}\left[H_{1}\ell_{1R}+H_{2}\ell_{2R}\right]-\frac{1}{\sqrt{2}}Y_{5}^{\ell}\left[\bar{L}_{1}H_{1}+\bar{L}_{2}H_{2}\right]\ell_{SR}+\text{ h.c. }
\end{align*}
and the mass matrix is given by
\begin{align}
M_{\ell}= & \begin{pmatrix}\mu_{1}^{\ell}+\mu_{2}^{\ell} & \sqrt{3}\mu_{2}^{\ell} & \sqrt{3}\mu_{5}^{\ell}\\
\sqrt{3}\mu_{2}^{\ell} & \mu_{1}^{\ell}-\mu_{2}^{\ell} & \mu_{5}^{\ell}\\
\sqrt{3}\mu_{4}^{\ell} & \mu_{4}^{\ell} & \mu_{3}^{\ell}
\end{pmatrix},\label{eq:ML_mass_matrix}
\end{align}
where\footnote{With $H_{i}=\frac{1}{\sqrt{2}}(\phi+i\phi')$ the neutral vacuum
expectation values are $\langle H_{i}^{0}\rangle=v_{i}/\sqrt{2}$, so each term of
$\mathcal{L}_{Y_{\ell}}$ contributes its explicit prefactor times $v_{i}/\sqrt{2}$;
this is what fixes the factors below. Note that $Y_{3}^{\ell}$ is the one structure
written without a $1/\sqrt{2}$ prefactor, hence the different factor in $\mu_{3}^{\ell}$;
adding a $1/\sqrt{2}$ there as well would make all five uniform.}
\begin{align*}
\mu_{1}^{\ell}= & \frac{v_{3}}{2}Y_{1}^{\ell} & \mu_{2}^{\ell}= & \frac{v_{2}}{2}Y_{2}^{\ell},\\
\mu_{3}^{\ell}= & \frac{v_{3}}{\sqrt{2}}Y_{3}^{\ell} & \mu_{4}^{\ell}= & \frac{v_{2}}{2}Y_{4}^{\ell},\\
 &  & \mu_{5}^{\ell}= & \frac{v_{2}}{2}Y_{5}^{\ell}.
\end{align*}
Matching the eigenvalues of (\ref{eq:ML_mass_matrix}) with the charged
lepton masses we obtain the following relations
\begin{align}
m_{e}= & \mu_{1}^{\ell}-2\mu_{2}^{\ell}\nonumber \\
m_{\mu}+m_{\tau}= & \mu_{1}^{\ell}+2\mu_{2}^{\ell}+\mu_{3}^{\ell}\nonumber \\
\left(m_{\tau}-m_{\mu}\right)^{2}= & {\displaystyle \left(2\mu_{3}^{\ell}-m_{\mu}-m_{\tau}\right)^{2}}+4\left(\mu_{4}^{\ell}+\mu_{5}^{\ell}\right)^{2}\label{eq:muil_equations}
\end{align}
where the last relation follows from $\left(\sigma_{1}-\sigma_{2}\right)^{2}=\text{tr}\left(A^{\top}A\right)-2\left|\det A\right|$
for the block $A$ defined in (\ref{eq:A_definition}) below, and is valid for
independent $\mu_{4}^{\ell}$ and $\mu_{5}^{\ell}$. (Only once $\mu_{4}^{\ell}=\mu_{5}^{\ell}$
has been established does it reduce to the more compact ${\displaystyle \left(2\mu_{3}^{\ell}-m_{\mu}-m_{\tau}\right)^{2}}+16\mu_{4}^{\ell}\mu_{5}^{\ell}$;
we keep the general form here because (\ref{eq:mu5_equal_mu4}) is derived from it.)
and we could solve as follows
\begin{align}
\mu_{1}^{\ell}= & \frac{1}{2}\left(m_{e}+m_{\mu}+m_{\tau}-\mu_{3}^{\ell}\right),\nonumber \\
\mu_{2}^{\ell}= & \frac{1}{4}\left(m_{\mu}+m_{\tau}-m_{e}-\mu_{3}^{\ell}\right),\label{eq:muil_sols}\\
\left(\mu_{4}^{\ell}+\mu_{5}^{\ell}\right)^{2}= & \left(m_{\tau}-\mu_{3}^{\ell}\right)\left(\mu_{3}^{\ell}-m_{\mu}\right),\nonumber
\end{align}
and
\begin{align*}
\mu_{1}^{\ell}+\mu_{2}^{\ell}= & \frac{1}{4}\left[m_{e}+3\left(m_{\mu}+m_{\tau}-\mu_{3}^{\ell}\right)\right],\\
\mu_{1}^{\ell}-\mu_{2}^{\ell}= & \frac{1}{4}\left[3m_{e}+\left(m_{\mu}+m_{\tau}-\mu_{3}^{\ell}\right)\right].
\end{align*}
Since the left-hand side of the last relation in (\ref{eq:muil_sols}) is a square,
its right-hand side must be non-negative, which imposes
\[
m_{\mu}\leq\mu_{3}^{\ell}\leq m_{\tau}
\]
with no further assumption.
In the case of $\mu_{5}^{\ell}=0$ $\left(\text{similarly to }\mu_{4}^{\ell}=0\right)$,
we have two cases:
\begin{enumerate}
\item $\mu_{3}^{\ell}=m_{\mu},$ then
\begin{align*}
\mu_{1}^{\ell}= & \frac{1}{2}\left(m_{e}+m_{\tau}\right)\\
\mu_{2}^{\ell}= & \frac{1}{4}\left(m_{\tau}-m_{e}\right)\\
\mu_{1}^{\ell}+\mu_{2}^{\ell}= & \frac{1}{4}\left(m_{e}+3m_{\tau}\right)\\
\mu_{1}^{\ell}-\mu_{2}^{\ell}= & \frac{1}{4}\left(3m_{e}+m_{\tau}\right)
\end{align*}
\[
M_{\ell}=\begin{pmatrix}\frac{1}{4}\left(m_{e}+3m_{\tau}\right) & \frac{\sqrt{3}}{4}\left(m_{\tau}-m_{e}\right) & 0\\
\frac{\sqrt{3}}{4}\left(m_{\tau}-m_{e}\right) & \frac{1}{4}\left(3m_{e}+m_{\tau}\right) & 0\\
\sqrt{3}\mu_{4}^{\ell} & \mu_{4}^{\ell} & m_{\mu}
\end{pmatrix}.
\]
\item $\mu_{3}^{\ell}=m_{\tau}$, so
\begin{align*}
\mu_{1}^{\ell}= & \frac{1}{2}\left(m_{e}+m_{\mu}\right)\\
\mu_{2}^{\ell}= & \frac{1}{4}\left(m_{\mu}-m_{e}\right)\\
\mu_{1}^{\ell}+\mu_{2}^{\ell}= & \frac{1}{4}\left(m_{e}+3m_{\mu}\right)\\
\mu_{1}^{\ell}-\mu_{2}^{\ell}= & \frac{1}{4}\left(3m_{e}+m_{\mu}\right)
\end{align*}
\[
M_{\ell}=\begin{pmatrix}\frac{1}{4}\left(m_{e}+3m_{\mu}\right) & \frac{\sqrt{3}}{4}\left(m_{\mu}-m_{e}\right) & 0\\
\frac{\sqrt{3}}{4}\left(m_{\mu}-m_{e}\right) & \frac{1}{4}\left(3m_{e}+m_{\mu}\right) & 0\\
\sqrt{3}\mu_{4}^{\ell} & \mu_{4}^{\ell} & m_{\tau}
\end{pmatrix}.
\]
\end{enumerate}

\subsection{Physical basis}

The diagonalization occurs in two stages, the first stage consists
of a rotation by $\pi/3$ $e-\mu$ sector, 
\begin{equation}
O_{12}=\begin{pmatrix}\frac{1}{2} & \frac{\sqrt{3}}{2} & 0\\
-\frac{\sqrt{3}}{2} & \frac{1}{2} & 0\\
0 & 0 & 1
\end{pmatrix}.\label{eq:O12_definition}
\end{equation}
The rotation $O_{12}$ gives a block diagonal matrix 
\begin{equation}
\tilde{M}_{\ell}=O_{12}^{\top}M_{\ell}O_{12}=\begin{pmatrix}m_{e} & 0\\
0 & A
\end{pmatrix}\label{eq:O12_rotation}
\end{equation}
where 
\begin{equation}
A=\begin{cases}
\begin{pmatrix}m_{\alpha} & 0\\
2\mu_{4}^{\ell} & m_{\beta}
\end{pmatrix} & \text{case 1 (\ensuremath{\alpha=\tau} and \ensuremath{\beta=\mu}) and case 2 (\ensuremath{\alpha=\mu} and \ensuremath{\beta=\tau}),}\\
\begin{pmatrix}\rho & 2\mu_{5}^{\ell}\\
2\mu_{4}^{\ell} & \mu_{3}^{\ell}
\end{pmatrix} & \text{in general},
\end{cases}\label{eq:A_definition}
\end{equation}
where
\[
\rho=m_{\mu}+m_{\tau}-\mu_{3}^{\ell}
\]
Then we need to diagonalize $A$, and then the lepton masses are equivalent
to the singular values of $A$, which arise from 
\[
N=A^{\dagger}A=A^{\top}A
\]
where the last identity is true with $\mu_{4}^{\ell}\mu_{5}^{\ell}\geq0$.
Explicitly
\[
N=\begin{cases}
\begin{pmatrix}4\left(\mu_{4}^{\ell}\right)^{2}+m_{\alpha}^{2} & 2\mu_{4}^{\ell}m_{\beta}\\
2\mu_{4}^{\ell}m_{\beta} & m_{\beta}^{2}
\end{pmatrix} & \text{case 1 (\ensuremath{\alpha=\tau} and \ensuremath{\beta=\mu}) and case 2 (\ensuremath{\alpha=\mu} and \ensuremath{\beta=\tau}),}\\
\begin{pmatrix}4\left(\mu_{4}^{\ell}\right)^{2}+\rho^{2} & 2\mu_{3}^{\ell}\mu_{4}^{\ell}+2\mu_{5}^{\ell}\rho\\
2\mu_{3}^{\ell}\mu_{4}^{\ell}+2\mu_{5}^{\ell}\rho & \left(\mu_{3}^{\ell}\right)^{2}+4\left(\mu_{5}^{\ell}\right)^{2}
\end{pmatrix} & \text{in general}.
\end{cases}
\]
comparing with the trace,
\begin{equation}
\text{tr}\left(N\right)=m_{\mu}^{2}+m_{\tau}^{2}\label{eq:TraceN}
\end{equation}
For cases 1 and 2
\[
\text{tr}\left(N\right)=4\left(\mu_{4}^{\ell}\right)^{2}+m_{\mu}^{2}+m_{\tau}^{2}
\]
consequently, the only one solution is 
\[
\mu_{4}^{\ell}=0,
\]
we conclude for these simplified cases that we only need $O_{12}$
to diagonalize $M_{\ell}$ and its diagonal form is $\tilde{M_{\ell}}$.
However, for the general case, from (\ref{eq:TraceN})
\begin{align}
\ensuremath{{\displaystyle 4\left[\left(\mu_{4}^{\ell}\right)^{2}+\left(\mu_{5}^{\ell}\right)^{2}\right]}} & =-2m_{\mu}m_{\tau}+2\mu_{3}^{\ell}\left(m_{\mu}+m_{\tau}\right)-2\left(\mu_{3}^{\ell}\right)^{2},\nonumber \\
 & =2\left(m_{\tau}-\mu_{3}^{\ell}\right)\left(\mu_{3}^{\ell}-m_{\mu}\right).\label{eq:trN_mu45}
\end{align}
Comparing with the third equation in (\ref{eq:muil_sols}), which involves
$\left(\mu_{4}^{\ell}+\mu_{5}^{\ell}\right)^{2}$ and is therefore independent
of (\ref{eq:trN_mu45}),
\begin{align*}
\left(\mu_{4}^{\ell}+\mu_{5}^{\ell}\right)^{2} & =\left(m_{\tau}-\mu_{3}^{\ell}\right)\left(\mu_{3}^{\ell}-m_{\mu}\right),\\
{\displaystyle 2\left[\left(\mu_{4}^{\ell}\right)^{2}+\left(\mu_{5}^{\ell}\right)^{2}\right]} & =\left(m_{\tau}-\mu_{3}^{\ell}\right)\left(\mu_{3}^{\ell}-m_{\mu}\right),
\end{align*}
so that subtracting the second from the first
\begin{align*}
\left(\mu_{4}^{\ell}+\mu_{5}^{\ell}\right)^{2}-2\left[\left(\mu_{4}^{\ell}\right)^{2}+\left(\mu_{5}^{\ell}\right)^{2}\right] & =0\\
-\left(\mu_{4}^{\ell}-\mu_{5}^{\ell}\right)^{2} & =0
\end{align*}
So, the only possibiity in this case is 
\begin{equation}
\mu_{5}^{\ell}=\mu_{4}^{\ell}.\label{eq:mu5_equal_mu4}
\end{equation}
In consequence, from (\ref{eq:mu5_equal_mu4}), $A$ is symmetric
and
\[
A=\begin{pmatrix}\rho & 2\mu_{4}^{\ell}\\
2\mu_{4}^{\ell} & \mu_{3}^{\ell}
\end{pmatrix},
\]
which is diagonalized by 
\[
o_{23}=\begin{pmatrix}\cos\left(\theta_{\ell}\right) & \sin\left(\theta_{\ell}\right)\\
-\sin\left(\theta_{\ell}\right) & \cos\left(\theta_{\ell}\right)
\end{pmatrix},
\]
with the angle given by
\begin{align*}
\tan\left(2\theta_{\ell}\right) & =\frac{4\mu_{4}^{\ell}}{2\mu_{3}^{\ell}-\left(m_{\mu}+m_{\tau}\right)}.
\end{align*}
Without loss of generality, we assume that $\mu_{4}^{\ell}>0$, then,
from (\ref{eq:muil_sols}) and (\ref{eq:mu5_equal_mu4})
\[
\mu_{4}^{\ell}=\frac{1}{2}\sqrt{\left(\mu_{3}^{\ell}-m_{\mu}\right)\left(m_{\tau}-\mu_{3}^{\ell}\right)}=\frac{1}{2}\sqrt{p_{1}p_{2}},\quad m_{\mu}<\mu_{3}^{\ell}<m_{\tau}.
\]
as a consequence
\[
\tan\left(2\theta_{\ell}\right)=\frac{2\sqrt{p_{1}p_{2}}}{p_{1}-p_{2}}
\]
where
\[
p_{1}=\mu_{3}^{\ell}-m_{\mu};\quad p_{2}=m_{\tau}-\mu_{3}^{\ell};\quad p_{1,2}>0.
\]
\begin{align*}
\sin\left(\theta_{\ell}\right)=\sqrt{\frac{p_{2}}{\Delta}} & \cos\left(\theta_{\ell}\right)=\sqrt{\frac{p_{1}}{\Delta}};\quad\Delta=p_{1}+p_{2}=m_{\tau}-m_{\mu}>0.
\end{align*}
 Finally, the full charged lepton mass matrix is diagonalized by an
orthogonal matrix given as follows 
\begin{equation}
O=O_{12}O_{23}\label{eq:O_definition}
\end{equation}
 where $O_{12}$ is given in (\ref{eq:O12_definition}) and 
\begin{equation}
O_{23}=\begin{pmatrix}1 & 0 & 0\\
0 & \cos\left(\theta_{\ell}\right) & \sin\left(\theta_{\ell}\right)\\
0 & -\sin\left(\theta_{\ell}\right) & \cos\left(\theta_{\ell}\right)
\end{pmatrix}.\label{eq:O23}
\end{equation}


\section{Lepton flavor violating Higgs decays}

From now on, we will consider the condition $\mu_{5}^{\ell}=\mu_{4}^{\ell}\left(Y_{5}^{\ell}=Y_{4}^{\ell}\right)$,
so, the charged lepton mass matrix is given by
\begin{align*}
M_{\ell}= & \begin{pmatrix}\mu_{1}^{\ell}+\mu_{2}^{\ell} & \sqrt{3}\mu_{2}^{\ell} & \sqrt{3}\mu_{4}^{\ell}\\
\sqrt{3}\mu_{2}^{\ell} & \mu_{1}^{\ell}-\mu_{2}^{\ell} & \mu_{4}^{\ell}\\
\sqrt{3}\mu_{4}^{\ell} & \mu_{4}^{\ell} & \mu_{3}^{\ell}
\end{pmatrix}
\end{align*}
which can be decomposed as follows
\[
M_{\ell}=v_{1}G_{1}+v_{2}G_{2}+v_{3}G_{S}
\]
where
\begin{align*}
G_{1}= & \frac{1}{v_{1}}\begin{pmatrix}0 & \sqrt{3}\mu_{2}^{\ell} & \sqrt{3}\mu_{4}^{\ell}\\
\sqrt{3}\mu_{2}^{\ell} & 0 & 0\\
\sqrt{3}\mu_{4}^{\ell} & 0 & 0
\end{pmatrix}=\frac{1}{v_{2}}\begin{pmatrix}0 & \mu_{2}^{\ell} & \mu_{4}^{\ell}\\
\mu_{2}^{\ell} & 0 & 0\\
\mu_{4}^{\ell} & 0 & 0
\end{pmatrix}\\
G_{2}= & \frac{1}{v_{2}}\begin{pmatrix}\mu_{2}^{\ell} & 0 & 0\\
0 & -\mu_{2}^{\ell} & \mu_{4}^{\ell}\\
0 & \mu_{4}^{\ell} & 0
\end{pmatrix}\\
G_{S}= & \frac{1}{v_{3}}\begin{pmatrix}\mu_{1}^{\ell} & 0 & 0\\
0 & \mu_{1}^{\ell} & 0\\
0 & 0 & \mu_{3}^{\ell}
\end{pmatrix}
\end{align*}
each $G_{k}$ couples with $H_{k}$ ($k=1,2,S$). So, in the lepton
mass basis, these matrices are transform by $O$ matrix , such as
follows
\begin{align}
\widetilde{G}_{k}= & O^{\top}G_{k}O.\label{eq:Gktilde}
\end{align}

In terms of $\widetilde{G}_{k}$ from (\ref{eq:Gktilde}), the charged
lepton interactions with neutral scalar are given as follows
\begin{equation}
\mathcal{L}_{\text{LFV}}^{h}=-h_{1}\overline{\ell}Q_{1}P_{R}\ell-h_{0}\overline{\ell}Q_{2}P_{R}\ell-h_{2}\overline{\ell}Q_{3}P_{R}\ell+\text{h.c.}\label{eq:L_LFVH}
\end{equation}
where
\begin{equation}
Q_{i}=\sum_{j=1,2,3}\left(R_{S}\right)_{ji}\widetilde{G}_{j}.\label{eq:Q_general}
\end{equation}


\subsection{Scenario A}

In this scenario, $R_{S}=R_{A}$, in consequence, 
\begin{align*}
Q_{1}\left(A\right)= & \frac{1}{v}\text{diag}\left(m_{e},m_{\mu},m_{\tau}\right),\\
Q_{2}\left(A\right)= & \begin{pmatrix}0 & -\frac{\sqrt{p_{1}}\left(m_{e}-\mu_{3}+p_{1}-3p_{2}\right)}{4v_{2}\sqrt{p_{1}+p_{2}}} & -\frac{\sqrt{p_{2}}\left(m_{e}-\mu_{3}+3p_{1}-p_{2}\right)}{4v_{2}\sqrt{p_{1}+p_{2}}}\\
-\frac{\sqrt{p_{1}}\left(m_{e}-\mu_{3}+p_{1}-3p_{2}\right)}{4v_{2}\sqrt{p_{1}+p_{2}}} & 0 & 0\\
-\frac{\sqrt{p_{2}}\left(m_{e}-\mu_{3}+3p_{1}-p_{2}\right)}{4v_{2}\sqrt{p_{1}+p_{2}}} & 0 & 0
\end{pmatrix}\\
Q_{3}\left(A\right)= & \begin{pmatrix}\frac{2m_{e}v_{2}}{vv_{3}}-\frac{v\left(m_{e}-\mu_{3}+p_{1}-p_{2}\right)}{4v_{2}v_{3}} & 0 & 0\\
0 & \frac{2m_{\mu}v_{2}}{vv_{3}}+\frac{p_{1}v\left(m_{e}-m_{\mu}+3p_{2}\right)}{4v_{2}v_{3}\left(p_{1}+p_{2}\right)} & -\frac{\sqrt{p_{1}}\sqrt{p_{2}}v\left(-m_{e}+\mu_{3}+p_{1}-p_{2}\right)}{4v_{2}v_{3}\left(p_{1}+p_{2}\right)}\\
0 & -\frac{\sqrt{p_{1}}\sqrt{p_{2}}v\left(-m_{e}+\mu_{3}+p_{1}-p_{2}\right)}{4v_{2}v_{3}\left(p_{1}+p_{2}\right)} & \frac{2m_{\tau}v_{2}}{vv_{3}}-\frac{p_{2}v\left(-m_{e}+m_{\tau}+3p_{1}\right)}{4v_{2}v_{3}\left(p_{1}+p_{2}\right)}
\end{pmatrix}
\end{align*}


\subsection{Scenario B}

Now, 

\begin{align*}
Q_{1}\left(B\right)= & \begin{pmatrix}-\frac{2m_{e}v_{2}}{vv_{3}}+\frac{v\left(m_{e}-m_{\mu}-p_{2}\right)}{4v_{2}v_{3}} & 0 & 0\\
0 & -\frac{2m_{\mu}v_{2}}{vv_{3}}-\frac{p_{1}v\left(m_{e}-m_{\mu}+3p_{2}\right)}{4v_{2}v_{3}\left(p_{1}+p_{2}\right)} & \frac{\sqrt{p_{1}}\sqrt{p_{2}}v\left(-m_{e}+\mu_{3}+p_{1}-p_{2}\right)}{4v_{2}v_{3}\left(p_{1}+p_{2}\right)}\\
0 & \frac{\sqrt{p_{1}}\sqrt{p_{2}}v\left(-m_{e}+\mu_{3}+p_{1}-p_{2}\right)}{4v_{2}v_{3}\left(p_{1}+p_{2}\right)} & -\frac{2m_{\tau}v_{2}}{vv_{3}}+\frac{p_{2}v\left(-m_{e}+m_{\tau}+3p_{1}\right)}{4v_{2}v_{3}\left(p_{1}+p_{2}\right)}
\end{pmatrix},\\
Q_{2}\left(B\right)= & \begin{pmatrix}0 & -\frac{\sqrt{p_{1}}\left(m_{e}-\mu_{3}+p_{1}-3p_{2}\right)}{4v_{2}\sqrt{p_{1}+p_{2}}} & -\frac{\sqrt{p_{2}}\left(m_{e}-\mu_{3}+3p_{1}-p_{2}\right)}{4v_{2}\sqrt{p_{1}+p_{2}}}\\
-\frac{\sqrt{p_{1}}\left(m_{e}-\mu_{3}+p_{1}-3p_{2}\right)}{4v_{2}\sqrt{p_{1}+p_{2}}} & 0 & 0\\
-\frac{\sqrt{p_{2}}\left(m_{e}-\mu_{3}+3p_{1}-p_{2}\right)}{4v_{2}\sqrt{p_{1}+p_{2}}} & 0 & 0
\end{pmatrix}\\
Q_{3}\left(B\right)= & \frac{1}{v}\text{diag}\left(m_{e},m_{\mu},m_{\tau}\right)
\end{align*}


\subsection{Scenario C}

But, 

\begin{align*}
Q_{1}\left(C\right)\approx & \frac{1}{v}\cos\left(\delta\right)\text{diag}\left(m_{e},m_{\mu},m_{\tau}\right)+\sin\left(\delta\right)Q_{1}\left(B\right),\\
Q_{2}\left(C\right)\approx & \begin{pmatrix}0 & -\frac{\sqrt{p_{1}}\left(m_{e}-\mu_{3}+p_{1}-3p_{2}\right)}{4v_{2}\sqrt{p_{1}+p_{2}}} & -\frac{\sqrt{p_{2}}\left(m_{e}-\mu_{3}+3p_{1}-p_{2}\right)}{4v_{2}\sqrt{p_{1}+p_{2}}}\\
-\frac{\sqrt{p_{1}}\left(m_{e}-\mu_{3}+p_{1}-3p_{2}\right)}{4v_{2}\sqrt{p_{1}+p_{2}}} & 0 & 0\\
-\frac{\sqrt{p_{2}}\left(m_{e}-\mu_{3}+3p_{1}-p_{2}\right)}{4v_{2}\sqrt{p_{1}+p_{2}}} & 0 & 0
\end{pmatrix}\\
Q_{3}\left(C\right)\approx & \cos\left(\delta\right)Q_{3}\left(A\right)+\frac{\sin\left(\delta\right)}{v}\text{diag}\left(m_{e},m_{\mu},m_{\tau}\right)
\end{align*}

\bibliographystyle{plain}
\bibliography{biblioteca}

\end{document}
